\section{Powers of finite sums}
\label{sec:finite_sum_powers}

The word expansion in the approximate ground-state projector of~\cite{Landau2013PolynomialTime} is valid without
commutation of the local operators. The following algebraic statement
preserves the order of every factor.

\begin{theorem}[Ordered expansion of a finite-sum power]
    \label{thm:finite_sum_power_ordered_words}
    \lean{Fintype.sum_pow_eq_sum_ordered_words}
    \leanok
    Let $R$ be a semiring, let $I$ be a finite set, and let $P_i\in R$
    for $i\in I$. For every integer $m\ge0$,
    \begin{align}
        \left(\sum_{i\in I}P_i\right)^m
        =\sum_{w\in I^{\{0,\ldots,m-1\}}}
            P_{w(0)}P_{w(1)}\cdots P_{w(m-1)}.
    \end{align}
    The factors on the right occur in the indicated order.
    For $m=0$, the unique empty word contributes the multiplicative identity,
    including when $I$ is empty. For $m>0$ and $I$ empty, both sides vanish.
\end{theorem}

\begin{proof}
    \leanok
    Induct on $m$. At $m=0$, the empty product gives the assertion.
    In the induction step, write
    $(\sum_i P_i)^{m+1}=(\sum_i P_i)(\sum_i P_i)^m$.
    Distribute multiplication over both sums. Every word of length $m+1$
    is uniquely determined by its first letter and its remaining word of
    length $m$, giving the required indexing of the sum.
\end{proof}

\begin{corollary}[Normalized complement expansion]
    \label{cor:normalized_complement_ordered_words}
    \lean{Fintype.smul_sub_average_pow_eq_sum_ordered_words}
    \leanok
    Let $R$ be a unital complex algebra, let $I$ be a nonempty finite set
    of cardinality $N$, and let $h_i\in R$. For $c\in\mathbb C$ and every
    integer $m\ge0$,
    \begin{align}
        \left[c\left(1-\frac1N\sum_{i\in I}h_i\right)\right]^m
        =\left(\frac cN\right)^m
          \sum_{w\in I^{\{0,\ldots,m-1\}}}
            (1-h_{w(0)})\cdots(1-h_{w(m-1)}).
    \end{align}
    Complex scalars act through the algebra structure. The scalar $c$
    may vanish, and the factors need not commute.
\end{corollary}

\begin{proof}
    \leanok
    \uses{thm:finite_sum_power_ordered_words}
    Write $S=\sum_i h_i$. Since $N>0$,
    \begin{align}
        \frac1N\sum_i(1-h_i)=1-\frac1N S.
    \end{align}
    Take the $m$-th power, move the central scalar outside the power,
    and apply Theorem~\ref{thm:finite_sum_power_ordered_words} to the
    complements $1-h_i$.
\end{proof}

For the chain Hamiltonian of~\cite{Landau2013PolynomialTime}, there are
$N-1$ local terms. Add the final term $h_N=0$, so that its complement is
$1$, and choose $c=(1-E_0/N)^{-1}$. This gives the ordered expansion used
there. Spectral bounds on this operator are separate assertions.
